
\documentclass[a4paper,11pt]{article}

\usepackage[margin=2.5cm]{geometry}
\usepackage[T1]{fontenc}
\usepackage[utf8]{inputenc}
\usepackage{listings}
\usepackage{hyperref}

\title{Practical Work 5: The Longest Path}
\author{Nguyen Trong Bach \\ Student ID: 23BI14057}
\date{\today}

\begin{document}

\maketitle

\section{Introduction}
This practical implements a small MapReduce style project called
\emph{The Longest Path}. The goal is to analyse a set of text files
where each line contains one full file path (for example, the output
of the command \texttt{find /}) and to find the longest path or paths.
Instead of using a full MapReduce framework, I design a simple
MapReduce-like program in~C++.

\section{Input and Output}
\begin{itemize}
  \item \textbf{Input}: a set of text files, each containing one full
  file path per line. Each file can be considered as the data coming
  from one laptop.
  \item \textbf{Output}: the maximum path length and all paths whose
  length is equal to that maximum.
\end{itemize}

\section{Mapper and Reducer Design}

\subsection*{Mapper}
Each input file is processed by a logical Mapper, implemented as the
function \texttt{map\_file(const std::string\&)}. For a given file the
Mapper:
\begin{enumerate}
  \item reads the file line by line;
  \item for each non-empty line computes the length of the path
  (number of characters);
  \item keeps track of the current maximum length inside that file;
  \item stores all paths whose length is equal to this local maximum.
\end{enumerate}
The Mapper returns a structure \texttt{MapResult} containing the local
maximum length and the list of paths that achieve this length.

\subsection*{Reducer}
The Reducer is implemented by the function
\texttt{reduce\_result(const MapResult\&, std::size\_t\&, std::vector<std::string>\&)}.
It maintains two global variables: the global maximum length and the
list of all paths that achieve this maximum. For each \texttt{MapResult}
received from a Mapper the Reducer:
\begin{itemize}
  \item replaces the global result if the local maximum is larger;
  \item appends the local paths if the local maximum is equal to the
        global maximum;
  \item ignores the local result if it is smaller.
\end{itemize}
This behaviour is analogous to MapReduce, where all partial results
for the same key are combined in the Reduce phase.

\section{System Organization}
The whole system is implemented in a single source file
\texttt{longest\_path.cpp}. The main program follows these steps:
\begin{enumerate}
  \item read the list of input file names from the command line;
  \item for each file, call the Mapper \texttt{map\_file} to obtain a
  \texttt{MapResult};
  \item call the Reducer \texttt{reduce\_result} to combine all
  \texttt{MapResult} objects into a global result;
  \item after all files have been processed, print the global maximum
  path length and the corresponding paths.
\end{enumerate}

\section{Implementation Snippet}
An example excerpt from the Mapper function is shown below:

\begin{lstlisting}[language=C++]
MapResult map_file(const std::string &filename) {
    MapResult result;
    std::ifstream in(filename);
    if (!in) {
        std::cerr << "Cannot open input file: " << filename << "\n";
        return result;
    }

    std::string line;
    while (std::getline(in, line)) {
        std::size_t len = line.size();
        if (len == 0) continue;

        if (len > result.max_len) {
            result.max_len = len;
            result.paths.clear();
            result.paths.push_back(line);
        } else if (len == result.max_len) {
            result.paths.push_back(line);
        }
    }
    return result;
}
\end{lstlisting}

\section{Group Work}
This practical work was completed individually. I was responsible for
designing the Mapper and Reducer logic, implementing the C++ program
and writing this report.

\section{Conclusion}
The Longest Path practical shows how the MapReduce idea can be applied
to a different problem than Word Count. By clearly separating the
Mapper and Reducer stages, the code is easy to understand and could
later be extended to run in parallel using multiple processes or
machines.

\end{document}
